\section{系统边界、统一口径与变量参数字典}
\label{sec:symbols}

\subsection{系统边界、执行链与责任划分}

本章空间边界从风、光、潮等发电单元的设备端资源输入开始，经海上集电与公共交流母线，
终止于四类价值交付点：电力海缆陆上受端、氢气管输或船运交付点、满足通信及服务等级
约束的算力交付点，以及海洋用能设备电气入口。源侧集电损耗、公共辅机、海缆损耗、
电解槽系统用电和数据中心PUE均由唯一责任模块计量，公共母线不重复扣减。

本章的运行优化以已建设设备容量为上限，不把风电、光伏、潮流能、海缆、电解槽、储氢
或算力容量作为投资决策变量。容量候选、资本金结构、融资约束和回收期属于第6章的投资
与产业落地评价；第4章仅为候选方案提供可复算的运行与价值结果。“风电占比85\%以上”
在本章作为项目定位或容量情景输入\assumptiontag，不是小时调度模型推导出的经济最优
装机比例。

模型执行链固定为“边界发布—联合决策—物理响应—母线总账—状态提交—价值评价”。
设备模块发布可用域和期初状态，4.9在同一可行域中形成调度决策，4.7只使用实际测点值
核账，4.8依据最终交付量评价。若接口联调保留“请求值—接受值—实际值”链路，该链路
用于表达设备响应，不在4.9之外形成第二套优化器。各模块唯一责任见表~\ref{tab:module-responsibility}。

\begin{table}[H]
\centering
\caption{第4章模块责任及主要接口}
\label{tab:module-responsibility}
\begin{tabularx}{\textwidth}{>{\centering\arraybackslash}p{1.2cm} p{4.2cm} X}
\toprule
小节 & 唯一责任 & 向联合模型提供的主要量 \\
\midrule
4.2 & 风、光、潮资源到公共汇集点的功率映射 & 逐源可用功率、状态、爬坡与场景边界 \\
4.3 & 电、氢、算力及海洋服务的交付通道 & 通道容量、损耗、接纳与最终交付量 \\
4.4 & 电池能量状态和构网备用代理 & 充放电功率、SOC、备用与期末状态 \\
4.5 & 电解制氢、储氢及可选氢转电 & 电解功率、产氢、库存、提取与回发边界 \\
4.6 & 任务--IT功率--设施功率--算力服务映射 & 设施功率、服务量、队列、PUE及SLA \\
4.7 & 公共母线唯一功率总账 & 平衡残差、实际注入、实际消耗与弃电分解 \\
4.8 & 经济--环境--可靠性目标和价值分账 & 目标向量、KPI及管理偏好接口 \\
4.9 & 全局物理、逻辑、服务与事件约束 & 联合可行域、比例规则与终端边界 \\
4.10 & 因果逐时求解、状态滚动和动态复核 & 调度动作、状态轨迹、求解及审计状态 \\
\bottomrule
\end{tabularx}
\sourcetype{项目内部模块责任定义。}
\end{table}

\subsection{公共测点、功率方向、损耗归属与统一时间尺度}

\subsubsection{集合、统一时间轴与信息边界}

令调度时段集合、源集合和价值路径集合分别为
\begin{equation}
\mathcal T=\{1,2,\ldots,T\},\qquad
\mathcal S=\{w,pv,tc\},\qquad
\mathcal J=\{E,H,C,M\},\qquad \Delta t_t>0.
\label{eq:common-sets}
\end{equation}
波浪能当前只保留接口，未取得二维功率矩阵前不进入主运行模型。主运行模型采用
$\Delta t_t=1\myunit{h}$\assumptiontag；功率为区间平均MW，电量为MWh，氢气产量与
库存为kg，算力服务功率为MW-IT，任务工作量为MWh-IT或MWh-CS。时段$t$可用信息集
$\mathcal I_t$只包含截至$t$的观测、当时已发布预警、合同参数和$t-1$末状态，不包含
$t+1$之后的真实出力、价格、负荷或全年总供电量。

\subsubsection{公共测点与功率正方向}

表~\ref{tab:meter-points}冻结公共测点。测点编码沿用V2机器接口，因此编码中的“03、04”
是接口版本的历史标识，不随本报告小节上位而改名；报告小节责任以表~\ref{tab:module-responsibility}
为准。

\begin{table}[H]
\centering
\footnotesize
\caption{公共测点、计量含义与责任模块}
\label{tab:meter-points}
\begin{tabularx}{\textwidth}{>{\ttfamily\raggedright\arraybackslash}p{4.0cm} X p{2.5cm}}
\toprule
测点 & 计量含义 & 责任模块 \\
\midrule
MP-DEVICE & 单台或单阵列设备电气端毛功率 & 4.2源侧 \\
MP-03-SOURCE-POI & 已扣源侧集电损耗的多源汇集功率 & 4.2源侧 \\
MP-04-COMMON-BUS & 所有源、储、算、用端口的公共功率守恒点 & 4.7总账 \\
MP-ELECTROLYZER & 电解制氢系统交流侧电气入口 & 4.5制氢 \\
MP-DC-FACILITY & 数据中心设施总用电入口 & 4.6绿算 \\
MP-CABLE-SEND & 海缆海上送端 & 4.3交付 \\
MP-CABLE-RECEIVE & 海缆陆上受端及电力结算点 & 4.3交付 \\
MP-H2-DELIVERY & 扣除物流损失后的氢产品交付点 & 4.3交付 \\
MP-COMPUTE-DELIVERY & 满足通信与SLA的算力服务交付点 & 4.3交付 \\
\bottomrule
\end{tabularx}
\sourcetype{项目内部冻结接口。}
\end{table}

优化层统一使用非负端口幅值：发电、储能放电和经允许的外购电为正注入；电池充电、
电解槽、数据中心、海洋负荷、海缆送端、公共辅机、损耗和弃电为正消耗。电池净功率只作
派生量：
\begin{equation}
P_t^{B,net}=P_t^{dis}-P_t^{ch},
\label{eq:bess-net-interface}
\end{equation}
其正值表示向公共母线放电。式~\eqref{eq:bess-net-interface}为项目接口定义，不替代4.4的
能量状态方程。每个可控端口保留$P^{req}\rightarrow P^{acc}\rightarrow P^{act}$链路，
母线平衡和价值评价只读取$P^{act}$或最终交付量。

\subsubsection{损耗与辅机唯一归属}

\begin{table}[H]
\centering
\footnotesize
\caption{损耗、辅机与产品损失的唯一归属}
\label{tab:loss-ownership}
\begin{tabularx}{\textwidth}{p{5.0cm} p{2.6cm} X}
\toprule
项目 & 唯一计算责任 & 公共总账处理 \\
\midrule
风机阵列、光伏场内电缆及潮流能集电损耗 & 4.2 & 在MP-03上游扣除，4.7不再扣除 \\
风、光、潮设备辅机 & 4.2计算、4.7记账 & 作为独立源侧辅机负荷扣一次 \\
公共升压、变换、母线和岛级公用系统 & 4.7 & 作为公共辅机或POI后损耗扣一次 \\
电池充放电损耗 & 4.4 & 只体现在能量状态方程，不在4.7另扣 \\
电解SEC、BOP、压缩纯化和储氢循环能耗 & 4.5 & 按已声明设备边界合成或单列，不重复包含 \\
IT、冷却、泵、网络和存储辅助功率 & 4.6 & 合成为设施总功率进入MP-DC-FACILITY \\
海缆损耗 & 4.3 & 4.7使用送端实际功率，收益使用受端电量 \\
管输或船运氢损失 & 4.3 & 只减少最终交付量，不回减4.5产氢量 \\
海洋用能设备电耗 & 4.3 & 实际服务功率进入4.7公共母线 \\
\bottomrule
\end{tabularx}
\sourcetype{项目内部能量与产品账簿规则。}
\end{table}

\subsubsection{统一时间尺度与聚合规则}

\begin{table}[H]
\centering
\caption{模型层级与统一时间尺度}
\label{tab:time-scales}
\begin{tabularx}{\textwidth}{p{3.3cm} p{2.6cm} X}
\toprule
层级 & 当前时间尺度 & 主要用途与限制 \\
\midrule
源侧原始数据/仿真 & 5 min & 风光潮波动、状态机和预测误差，进入主模型前聚合 \\
主运营调度 & 1 h & 48 h机制验证与8,760 h因果滚动，不表征秒级动态 \\
快速可行性复核 & 5--15 min & 源侧快速爬坡、PCS和电解槽动态边界复核 \\
构网正序/EMT/HIL & 不大于1 s & 频率、电压、限流、故障穿越和黑启动，独立于小时MILP \\
规划与财务 & 1 year & 建设、替换、折旧、现金流和生命周期评价 \\
\bottomrule
\end{tabularx}
\sourcetype{项目建模假设；具体步长待场址数据、控制器和仿真平台校准。}
\end{table}

从细时间尺度聚合至小时尺度时，功率按能量守恒求平均，电量、产氢量和服务量求和，SOC、
氢库存、任务队列和设备状态取时段末值；故障、保护和复归事件不作线性插值，极值与爬坡
指标另行保留。小时模型不因聚合而获得对未来真实轨迹的预见。

\subsection{全模型变量与参数字典}

\subsubsection{变量、状态与审计字段}

表~\ref{tab:variable-dictionary}按当前报告责任重新汇总原4.2冻结字典的联合模型字段。
机器实现仍以V2 Schema为唯一接口；本表用于报告解释，不建立第二套字段定义。$req$、
$acc$、$act$分别表示请求、设备接受和实际测量/执行值，逐时字段均携带时间戳、测点、
参数版本和质量标志。

\begingroup
\footnotesize
\begin{longtable}{p{1.2cm} p{3.45cm} p{1.55cm} p{5.0cm} p{2.15cm}}
\caption{联合模型变量、状态与审计字段字典}\label{tab:variable-dictionary}\\
\toprule
模块 & 符号/字段 & 单位 & 含义与域 & 生产责任 \\
\midrule
\endfirsthead
\multicolumn{5}{l}{\small 续表~\thetable}\\
\toprule
模块 & 符号/字段 & 单位 & 含义与域 & 生产责任 \\
\midrule
\endhead
公共 & $t,\Delta t_t,\omega,\pi_\omega$ & h/p.u. & 时段起点、步长、场景及概率，$\pi_\omega\ge0$且求和为1 & 公共数据层 \\
公共 & $A_{j,t},z_t^{evt}$ & p.u./枚举 & 设备可用/降额系数及正常、预警、过境、恢复事件状态 & 各设备/总控 \\
公共 & $P_{j,t}^{req/acc/act}$ & MW & 可控端口请求、接受和实际值，均采用非负幅值 & 4.9/设备模块 \\
4.2 & $P_{i,t}^{av,g},P_{i,t}^{loss,col}$ & MW & 源$i$设备端毛可用功率及集电损耗 & 源侧模块 \\
4.2 & $P_{i,t}^{av},P_{i,t}^{use},P_{i,t}^{curt}$ & MW & 公共汇集点可用、实际利用和调度弃电功率 & 源侧/4.9 \\
4.2 & $F_{i,t}^{up}$、$F_{i,t}^{down}$、$Q_{i,t}^{min/max}$ & MW；Mvar & 调节能力和无功边界，不等同已承诺备用 & 源侧模块 \\
4.3 & $P_t^{cab,s},P_t^{cab,r},P_t^{cab,loss}$ & MW & 海缆送端、陆上受端及通道损耗 & 交付模块 \\
4.3 & $m_t^{pipe},m_t^{ship},m_t^{H2,del}$ & kg/时段 & 管输提取、船运装载和最终氢交付量 & 交付模块 \\
4.3 & $W^{C,del},W^{SLA,fail}$ & MWh-CS & 可计费算力交付量及未满足合同的工作量 & 交付模块 \\
4.3 & $P_t^{mar,dem}$、$P_t^{mar}$、$P_t^{ENS,mar}$ & MW & 海洋用能需求、实际服务和未供功率 & 交付模块 \\
4.4 & $P_t^{ch}$、$P_t^{dis}$、$E_t^B$、$SOC_t$ & MW；MWh；p.u. & 电池充放电功率、时段末能量和荷电状态 & 电池模块 \\
4.4 & $u_t^{ch},u_t^{dis},g_t$ & 0--1 & 充放电互斥状态及构网责任状态 & 电池/4.9 \\
4.5 & $P_t^{EL},n_t^{EL},n_t^{start}$ & MW/个 & 电解系统功率、在线模块数和启动模块数 & 制氢模块 \\
4.5 & $m_t^{prod},V_t^{water},FLEH_t$ & kg/m$^3$/h & 时段产氢量、工艺给水量和累计等效满负荷小时 & 制氢模块 \\
4.5 & $H_t,m_t^{vent}$ & kg & 时段末储氢库存及独立记录的放散量 & 储氢模块 \\
4.5 & $P_t^{H2P}$、$m_t^{H2P}$、$u_t^{H2P}$ & MW；kg；0--1 & 可选氢转电功率、耗氢量和启用状态 & 韧性支路/4.9 \\
4.6 & $A_{k,t}^{C},x_{k,t},q_{k,t}$ & MWh-CS & 任务到达、完成量和时段末队列 & 算力模块 \\
4.6 & $P_t^{IT},P_t^{cool},P_t^{aux},P_t^{DC}$ & MW & IT、冷却、其他辅助和设施总功率 & 算力模块 \\
4.6 & $U_t^{C,r},U_t^{C,f},U_t^{C,s}$ & MWh-CS & 刚性任务未服务、柔性逾期和可选任务未接纳量 & 算力模块 \\
4.7 & $P_t^{int},P_t^{ENS,int},P_t^{spill}$ & MW & 岛内关键负荷、未供功率和显式弃电 & 4.9/总账 \\
4.7 & $r_t^P$ & MW & 公共母线平衡残差，闭合后应在容差内 & 总账模块 \\
4.8 & $F_{eco},F_{env},ENS,\eta_{RE}$ & 复合单位 & 经济、环境、可靠性目标和可再生能源利用率 & 评价模块 \\
4.9--4.10 & $\bm x_t,\bm y_t,\bm u_t$ & -- & 连续变量、整数变量和本时段提交动作 & 联合求解器 \\
审计 & schemaId、parameterVersion、qualityFlag & 文本/逻辑 & 接口、参数版本与数据有效性；无效数据不得静默填零 & 全部模块 \\
审计 & dispatchStatus、constraintCode & 枚举 & 求解状态及请求与实际值不一致的可追溯原因 & 4.10/设备模块 \\
\bottomrule
\end{longtable}
\endgroup
\sourcetype{项目V2冻结接口及V5联合模型字段；报告编号已按本版小节上位映射。}

\subsubsection{分模块参数及证据状态}

表~\ref{tab:module-parameters}把原分散在各子模型中的参数集中到4.1。参数表只负责名称、
单位、当前状态和校准来源；具体物理作用仍由4.2--4.10相应公式说明。除表内明确给出的
算力场景代理值外，工程数值一律从版本化场景文件读取，不在正文复制默认值。算力代理值
沿用现行V5模型配置，全部为\assumptiontag，并非设备商保证值。

\begingroup
\footnotesize
\begin{longtable}{p{1.15cm} p{3.25cm} p{1.65cm} p{3.15cm} p{4.15cm}}
\caption{第4章分模块参数字典与校准来源}\label{tab:module-parameters}\\
\toprule
模块 & 参数 & 单位 & 当前值或状态 & 首选来源类型/校准责任 \\
\midrule
\endfirsthead
\multicolumn{5}{l}{\small 续表~\thetable}\\
\toprule
模块 & 参数 & 单位 & 当前值或状态 & 首选来源类型/校准责任 \\
\midrule
\endhead
公共 & $\Delta t,\epsilon^P$ & h/MW & 1 h；平衡容差由求解精度配置\assumptiontag & 项目模型配置；仿真平台校准 \\
公共 & $\pi_\omega,A_{j,t},z_t^{evt}$ & p.u./枚举 & 版本化逐时输入\assumptiontag & 气象海洋数据、OEM/EPC、应急预案 \\
4.2 & $P_i^r,f_{P-v},f_{flood/ebb}$ & MW/曲线 & 场景容量与设备曲线输入\assumptiontag & OEM/型式试验、国家及国际标准 \\
4.2 & $R_i^{up/down}$、$Q_i^{min/max}$ & MW/h；Mvar & 逐设备输入\assumptiontag & OEM控制规范、并网与现场试验 \\
4.2 & $\rho_g,\gamma_P,k_{wake}$及状态阈值 & 复合单位 & 场址/设备输入\assumptiontag & 海洋实测、OEM、CFD与型式试验 \\
4.3 & $C^{cab},p_0^{cab},\alpha^{cab},\beta^{cab}$ & MW/复合单位 & 拓扑和长度相关输入\assumptiontag & 设计院、海缆/换流设备厂家 \\
4.3 & $\bar m^{pipe/ship}$、$\eta^{pipe/ship}$、$z^{window}$ & kg/时段；p.u. & 储运情景输入\assumptiontag & 储运企业、港航单位、承购合同 \\
4.3 & 光纤容量、时延及SLA阈值 & 复合单位 & 合同情景输入\assumptiontag & 通信运营商、客户合同、企业访谈 \\
4.4 & $E^{B,r}$、$P^{ch/dis,max}$、$SOC_{min/max}$ & MWh；MW；p.u. & 容量情景输入\assumptiontag & OEM规格、EPC设计与验收试验 \\
4.4 & $\eta_{ch/dis},\rho_B$ & p.u. & 设备曲线输入\assumptiontag & OEM效率/自放电试验、核心期刊论文 \\
4.4 & $P^{GFM,up/down}$、$E^{GFM,up/down}$ & MW；MWh & 构网备用输入\assumptiontag & PCS厂商、系统研究、REGFM/EMT/HIL \\
4.5 & $N^{EL}$、$P^{mod,min/max}$、$R^{EL,up/down}$ & 个；MW & 模块化设备输入\assumptiontag & 电解槽OEM及BOP集成商 \\
4.5 & $SEC,e^{start},w_{water}$ & kWh/kg等 & 工艺包输入\assumptiontag & OEM公开资料、政府目标、实测校准 \\
4.5 & $H^{max},\rho_H,H^{target}$ & kg/p.u. & 储氢与安全情景输入\assumptiontag & 储氢设备企业、安全评估与EPC \\
4.5 & $\eta^{H2P},P^{H2P,max}$ & p.u./MW & 默认资产关闭\assumptiontag & 待氢转电技术选型和OEM校准 \\
4.6 & $C^{IT}$ & MW & 7.08\assumptiontag & 现行V5算力场景配置；待企业校准 \\
4.6 & $r^{idle},\eta^{dist}$ & p.u. & 0.30；0.96\assumptiontag & 现行V5算力场景配置；待设备实测 \\
4.6 & $T^{ref},PUE^{design}$ & $^{\circ}$C/p.u. & 23.5；1.10\assumptiontag & 项目PUE锚点；PUE定义见行业协会/政府资料 \\
4.6 & $P_0^{cool},P^{aux}$ & MW & 0.05；0.10\assumptiontag & 现行V5算力场景配置；待液冷/网络设备校准 \\
4.6 & $k_1,k_2,k_T$ & 复合单位 & 0.03410556；0.0001；0.002\assumptiontag & 半机理代理反推；待实测曲线校准 \\
4.6 & $L$、$p_{GPU}^{eq}$ & h；kW/GPU & 12；0.90\assumptiontag & 任务SLA与经济换算代理；待客户合同校准 \\
4.8 & 电、氢、算力及海洋服务价格 & CNY/交付单位 & 分时合同输入\assumptiontag & 承购合同、市场规则、企业访谈一手数据 \\
4.8 & 排放因子、寿命与更换参数 & 复合单位 & 生命周期清单输入\assumptiontag & 政府/行业数据库、OEM、核心期刊论文 \\
4.8 & 折现率、评价期、币值基年 & p.u./year & 财务情景输入\assumptiontag & 项目融资方案、审计口径与第6章 \\
4.8--4.10 & $\varepsilon_{ENS},\varepsilon_{GHG},\varepsilon_{RE},\delta$ & 复合单位 & 管理阈值输入\assumptiontag & 项目决策者、合同底线与敏感性分析 \\
\bottomrule
\end{longtable}
\endgroup
\sourcetype{项目V5模型配置、国际/国家标准、政府官网、行业协会资料、制造商公开资料及核心期刊论文；未签认工程值均标注为假设值。}

\subsection{状态传递、数据质量与证据等级}

电池能量、氢库存、电解槽在线模块数及上一步功率、算力任务队列和设备在线状态均由其
责任模块在时段末提交，并原样成为下一时段初值。任何缺失初态、非有限值、单位错配、
测点错配或时间轴不连续均应在优化前报错，不得用模块默认值静默填充。运行输出至少包含
参数版本、时间戳、来源类型、质量标志、求解状态及母线/状态递推残差，以保证第5章逐时
账簿可追溯。

本章证据分为三类：国际/国家标准、政府或核心期刊文献用于支撑公式结构和评价口径；
OEM、EPC、承购合同及企业访谈用于校准项目参数；项目接口定义用于统一计量和软件契约。
尚未核验原始来源的专业代理公式标注为\verifytag，未获得场址、OEM、EPC、合同或实测
签认的具体参数标注为\assumptiontag。接口定义不冒充自然规律，工程参数也不以公开行业
均值替代项目签认值。
